%-----------------------------------------------------------------------
\subsection{Exergy analysis}
\label{sec:exergy}
%-----------------------------------------------------------------------

At cyclic steady state the store returns every unit of energy it absorbs,
Eq.~\eqref{eq:css-unity}, so the first law cannot express what the store, or
any other component, costs. The cost of each component is the exergy it
destroys, and the analysis below allocates the exergy input of the plant over
one complete cycle to the net electrical output, to a discarded stream and to
the destruction in each component.

\subsubsection{Dead state and reservoirs}

The dead state is taken as the ORC heat sink, $T_0=T_\mathrm{sink}$, at the
pressure of each stream. It is the lowest temperature in the system, so all
stream exergies are non-negative, and the heat rejected in the ORC condenser
carries no exergy. The plant exchanges heat with two reservoirs: the sink, at
the dead state, and the geothermal stream that supplies the HTHP evaporator.
The latter is not at the dead state and is an exergy input to the plant. The
exergy given up by a stream of water cooling from $T_a$ to $T_b$ at
$p_\mathrm{w}$ is
\begin{equation}
  \Exd_\mathrm{w}(T_a,T_b)=\md\left[h_\mathrm{w}(T_a)-h_\mathrm{w}(T_b)
     -T_0\left(s_\mathrm{w}(T_a)-s_\mathrm{w}(T_b)\right)\right],
  \label{eq:ex-stream}
\end{equation}
and $\Exd_\mathrm{w}(T_a,T_0)$ is the exergy of the stream relative to the
dead state.

\subsubsection{The geothermal source as a finite resource}
\label{sec:geo}

The heat-pump source is co-produced or geothermal water drawn from producing
wells at $T_\mathrm{source}$ and reinjected at
$T_\mathrm{4a}=T_\mathrm{source}-\Delta T_\mathrm{3A,4A}$. Its flow rate is not
prescribed but derived from the evaporator duty, in the same way as $\Nw$ is
derived from the storage duty. With
$\md_\mathrm{ref}=\dot Q_\mathrm{out,HP}/(h_\mathrm{2h}-h_\mathrm{3h})$ the
refrigerant flow through the condenser,
\begin{gather}
  \dot Q_\mathrm{src}=\md_\mathrm{ref}\left(1-x_\mathrm{8h}\right)
     \left(h_\mathrm{13h}-h_\mathrm{4h}\right),
  \notag\\
  \md_\mathrm{geo}=\frac{\dot Q_\mathrm{src}}
     {h_\mathrm{w}(T_\mathrm{source})-h_\mathrm{w}(T_\mathrm{4a})},
  \qquad
  N_\mathrm{geo}=\frac{\md_\mathrm{geo}}{\md_\mathrm{geo,well}},
  \label{eq:mgeo}
\end{gather}
where $\md_\mathrm{geo,well}$ is the yield of one producing well. Deriving
$\md_\mathrm{geo}$ places the uncertainty in the per-well yield, which is a
field-measured quantity, rather than in a plant-level flow rate; since the
source exergy is homogeneous of degree one in $\md_\mathrm{geo}$, the choice
does not affect the exergetic efficiencies below. $\Delta T_\mathrm{3A,4A}$ is
bounded above by a minimum reinjection temperature
$T_\mathrm{reinj,min}$, which represents formation constraints such as scaling
and injectivity.

Two conventions are possible for the exergy charged to the plant for the
geothermal stream. Under the \emph{resource} convention adopted here, the plant
is charged with the full exergy of the produced stream relative to the dead
state, and the exergy still carried by the reinjected stream is booked as a
loss attributed to no component:
\begin{gather}
  \Exd_\mathrm{geo}=\Exd_\mathrm{w}\!\left(T_\mathrm{source},T_0\right),
  \qquad
  \Exd_\mathrm{reinj}=\Exd_\mathrm{w}\!\left(T_\mathrm{4a},T_0\right),
  \notag\\
  u_\mathrm{geo}=1-\frac{\Exd_\mathrm{reinj}}{\Exd_\mathrm{geo}},
  \label{eq:ex-geo}
\end{gather}
where both terms are evaluated with $\md_\mathrm{geo}$ and $u_\mathrm{geo}$ is
the fraction of the resource that the plant uses. Under the alternative
\emph{stripped} convention, only the exergy removed in the evaporator,
$\Exd_\mathrm{geo}-\Exd_\mathrm{reinj}=\Exd_\mathrm{w}(T_\mathrm{source},
T_\mathrm{4a})$, is charged and reinjection is free. The resource convention is
appropriate for a dedicated producer, which is drilled, completed and operated
whether or not the exergy it lifts is used. The distinction is not cosmetic:
the two conventions give opposite guidance on $\Delta T_\mathrm{3A,4A}$
(Section~\ref{sec:res-exergy}). Under the stripped convention the exergetic
efficiency decreases monotonically with $\Delta T_\mathrm{3A,4A}$, which drives
the design toward the smallest cooling of the source and hence the largest
geothermal flow and well count; under the resource convention it has an
interior maximum.

\subsubsection{Exergy balance over a cycle}

Over one cycle at CSS the exergy of the store returns to its initial value, and
the balance for the whole plant reads
\begin{equation}
  W_\mathrm{el,in}+\Ex_\mathrm{geo}
  =W_\mathrm{el,out}+\Ex_\mathrm{reinj}+\sum_{j}I_{j},
  \label{eq:ex-balance}
\end{equation}
with $W_\mathrm{el,in}=\dot W_\mathrm{el,in}t_\mathrm{ch}$,
$W_\mathrm{el,out}=\dot W_\mathrm{el,out}t_\mathrm{dc}$,
$\Ex_\mathrm{geo}=\Exd_\mathrm{geo}t_\mathrm{ch}$ and
$\Ex_\mathrm{reinj}=\Exd_\mathrm{reinj}t_\mathrm{ch}$, all evaluated on the
energy the field actually moves at cyclic steady state rather than on the
target duties. The
destruction in component $j$ follows from the Gouy--Stodola theorem,
$I_j=T_0\,S_{\mathrm{gen},j}$, with the entropy generation computed from the
state points and flow fractions of that component alone.

\paragraph{Heat pump} Per unit mass through the condenser, and with the water
flow rates referred to $\md_\mathrm{ref}$,
\begin{align}
  s_{\mathrm{gen,C,HP}}&=s_\mathrm{2h}-s_\mathrm{6h},\notag\\
  s_{\mathrm{gen,C,LP}}&=\left(1-x_\mathrm{8h}\right)\left(s_\mathrm{5h}-s_\mathrm{1h}\right),
  \notag\\
  s_{\mathrm{gen,V,7-8}}&=s_\mathrm{8h}-s_\mathrm{7h},\notag\\
  s_{\mathrm{gen,V,9-4}}&=\left(1-x_\mathrm{8h}\right)\left(s_\mathrm{4h}-s_\mathrm{9h}\right),
  \notag\\
  s_{\mathrm{gen,sep}}&=x_\mathrm{8h}s_\mathrm{10h}+\left(1-x_\mathrm{8h}\right)s_\mathrm{9h}-s_\mathrm{8h},\notag\\
  s_{\mathrm{gen,R1}}&=\left(s_\mathrm{12h}-s_\mathrm{3h}\right)+x_\mathrm{8h}\left(s_\mathrm{11h}-s_\mathrm{10h}\right),
  \notag\\
  s_{\mathrm{gen,R2}}&=\left(s_\mathrm{7h}-s_\mathrm{12h}\right)+\left(1-x_\mathrm{8h}\right)\left(s_\mathrm{1h}-s_\mathrm{13h}\right),\notag\\
  s_{\mathrm{gen,cond}}&=\left(s_\mathrm{3h}-s_\mathrm{2h}\right)
     +\frac{\md_\mathrm{w,ch}^\mathrm{tot}}{\md_\mathrm{ref}}
      \left[s_\mathrm{w}(T_\mathrm{3c})-s_\mathrm{w}(T_\mathrm{2c})\right],
  \notag\\
  s_{\mathrm{gen,evap}}&=\left(1-x_\mathrm{8h}\right)\left(s_\mathrm{13h}-s_\mathrm{4h}\right)
     +\frac{\md_\mathrm{geo}}{\md_\mathrm{ref}}
      \left[s_\mathrm{w}(T_\mathrm{4a})-s_\mathrm{w}(T_\mathrm{source})\right],
  \label{eq:sgen-hp}
\end{align}
where $\md_\mathrm{w,ch}^\mathrm{tot}=\Nw n_t\md_\mathrm{ch}$ is the total
water flow, and $I_j=T_0\md_\mathrm{ref}s_{\mathrm{gen},j}t_\mathrm{ch}$. The
motor loss is
$I_\mathrm{mot}=(\dot W_\mathrm{el,in}-\dot W_\mathrm{C})t_\mathrm{ch}$, with
$\dot W_\mathrm{C}=\md_\mathrm{ref}[(h_\mathrm{2h}-h_\mathrm{6h})
+(1-x_\mathrm{8h})(h_\mathrm{5h}-h_\mathrm{1h})]$ the shaft power of the two
compressors.

\paragraph{Organic Rankine cycle} Per unit mass through the high-pressure
turbine, whose flow rate is
$\md_\mathrm{orc}=\dot Q_\mathrm{in,ORC}/[(h_\mathrm{3e}-h_\mathrm{10e})
+y(h_\mathrm{6e}-h_\mathrm{5e})]$,
\begin{align}
  s_{\mathrm{gen,T1}}&=s_\mathrm{5e}-s_\mathrm{3e},\notag\\
  s_{\mathrm{gen,T2}}&=y\left(s_\mathrm{4e}-s_\mathrm{6e}\right),
  \notag\\
  s_{\mathrm{gen,P1}}&=y\left(s_\mathrm{2e}-s_\mathrm{1e}\right),\notag\\
  s_{\mathrm{gen,P2}}&=\left(1-y\right)\left(s_\mathrm{8e}-s_\mathrm{7e}\right),
  \notag\\
  s_{\mathrm{gen,FH}}&=\left(1-y\right)\left(s_\mathrm{7e}-s_\mathrm{5e}\right)
     +y\left(s_\mathrm{9e}-s_\mathrm{2e}\right),\notag\\
  s_{\mathrm{gen,mix}}&=s_\mathrm{10e}-y\,s_\mathrm{9e}-\left(1-y\right)s_\mathrm{8e},
  \notag\\
  s_{\mathrm{gen,evap}}&=\left(s_\mathrm{3e}-s_\mathrm{10e}\right)
     +y\left(s_\mathrm{6e}-s_\mathrm{5e}\right)
     +\frac{\md_\mathrm{w,dc}^\mathrm{tot}}{\md_\mathrm{orc}}
      \left[s_\mathrm{w}(T_\mathrm{3d})-s_\mathrm{w}(T_\mathrm{2d})\right],\notag\\
  s_{\mathrm{gen,cond}}&=y\left(s_\mathrm{1e}-s_\mathrm{4e}\right)
     +\frac{y\left(h_\mathrm{4e}-h_\mathrm{1e}\right)}{T_0},
  \label{eq:sgen-orc}
\end{align}
with $I_j=T_0\md_\mathrm{orc}s_{\mathrm{gen},j}t_\mathrm{dc}$. Because the sink
is at the dead state, the whole exergy of the heat rejected in the condenser is
destroyed there. The generator loss is
$I_\mathrm{gen}=(\dot W_\mathrm{T}-\dot W_\mathrm{el,out})t_\mathrm{dc}$, with
$\dot W_\mathrm{T}$ the net shaft power of the turbines less the pumps.

\paragraph{Borehole field} The destruction in the borehole field is obtained as
the difference between the exergy the water gives up during charging and the
exergy it recovers during discharging,
\begin{equation}
  I_\mathrm{BB}=\Exd_\mathrm{w}^\mathrm{tot}\!\left(T_\mathrm{3c},T_\mathrm{2c}\right)t_\mathrm{ch}
   -\Exd_\mathrm{w}^\mathrm{tot}\!\left(T_\mathrm{2d},T_\mathrm{3d}\right)t_\mathrm{dc},
  \label{eq:I-BB}
\end{equation}
evaluated with the total water flow of each half-cycle. The expression is exact
at CSS, because the exergy of the PCM returns to its initial value over the
cycle and no entropy function of the PCM is required. It aggregates the
destruction over the depth and duration of the cycle and does not locate it;
resolving it in $s$ and $t$ requires an entropy relation $s(E')$ consistent
with Eq.~\eqref{eq:enth-curve}. \todo{independent evaluation of $I_\mathrm{BB}$
from $s(E')$ integrated over the march.}

\subsubsection{Exergetic efficiencies}

The round-trip efficiency, Eq.~\eqref{eq:rte}, books the geothermal source as
free. The exergetic efficiency under the resource convention charges it, and
the stripped value is reported alongside for comparison:
\begin{equation}
  \psi=\frac{W_\mathrm{el,out}-W_\mathrm{f,dc}}
            {W_\mathrm{el,in}+W_\mathrm{f,ch}+\Ex_\mathrm{geo}},
  \qquad
  \psi_\mathrm{str}=\frac{W_\mathrm{el,out}-W_\mathrm{f,dc}}
            {W_\mathrm{el,in}+W_\mathrm{f,ch}+\Ex_\mathrm{geo}-\Ex_\mathrm{reinj}},
  \label{eq:psi}
\end{equation}
where $W_\mathrm{f}=\dot W_\mathrm{f}t$ are the pumping energies of the
borehole field, so that $\psi$ and RTE share the same boundary and differ only
in the charge for the geothermal stream. The geothermal well count
$N_\mathrm{geo}$, Eq.~\eqref{eq:mgeo}, is reported together with $\Nw$, and the
total number of wells per unit of net output,
$(\Nw+N_\mathrm{geo})/\dot W_\mathrm{el,out}$, is carried as a design indicator:
a comparison based on RTE alone favours designs that buy COP with an
arbitrarily large geothermal flow.

Including the pumping terms is what makes $\psi$ and RTE share a control
volume, and it is not a small correction: in the reference case $\psi$ is
\num{0.246} if they are omitted and \num{0.236} if they are charged. The code
was brought onto this definition at v0.15, and a build-time check now asserts
that the RTE computed inside the audit and the RTE computed from the marched
field agree to the last bit --- a cross-check that exists because this
discrepancy was found while writing the present section rather than by any
test in the code.
